\section{Method}
\label{sec:method}

Our goal is to take $n_{\text{train}}$ samples in $d$ dimensions, form the
empirical SD-KDE (score + shift), and evaluate the resulting density on
$n_{\text{test}}$ queries.  Writing the computation in matrix form highlights
the parallel structure.  Let $x_i, y_j \in \mathbb{R}^d$ denote training and
query points, with bandwidth $h$ and squared Euclidean distance
\[
  \lVert x_i - y_j \rVert^2
  =
  \lVert x_i \rVert^2
  +
  \lVert y_j \rVert^2
  - 2\,x_i^\top y_j.
\]
Stacking training samples into $X \in \mathbb{R}^{n_{\text{train}}\times d}$
and queries into $Y \in \mathbb{R}^{n_{\text{test}}\times d}$, the
pairwise dot-product matrix
\[
  G
  =
  X Y^\top
  \in
  \mathbb{R}^{n_{\text{train}}\times n_{\text{test}}}
\]
dominates the arithmetic once $d$ is moderately large ($d\gtrsim 16$).  On
modern NVIDIA GPUs this GEMM can be mapped to Tensor Cores and evaluated at
5--10$\times$ the throughput of standard FP32 SIMT arithmetic, particularly
when we restrict to $d$ that are multiples of 16 and use Triton's
\texttt{tl.dot} interface.  The remaining operations---vector norms,
broadcasted additions, and exponentials---are all $O(n_{\text{train}}
 n_{\text{test}})$ but quickly become a small fraction of the total FLOPs as
$d$ grows.

The empirical SD-KDE score inherits the same GEMM structure.  Its numerator
involves terms of the form
\[
  \sum_j -(x_i - y_j)\,\varphi_{ij},
  \qquad
  \varphi_{ij}
  = \exp\!\left(-\frac{\lVert x_i - y_j\rVert^2}{2h^2}\right),
\]
which naively suggests $O(n_{\text{train}} n_{\text{test}} d)$ additional
elementwise arithmetic.  Using the identity
\[
  \sum_j (x_i - y_j)\,\varphi_{ij}
  =
  x_i \sum_j \varphi_{ij}
  -
  \sum_j \varphi_{ij}\,y_j,
\]
we can decompose the numerator into a second GEMM:
\[
  T
  =
  \Phi Y.
\]
Thus both the KDE evaluation and the SD-KDE score numerator reduce to
Tensor-Core-accelerable matrix multiplies, plus $O(n^2)$ scalar work for the
norms and exponentials.  In this note we focus on the $d=16$ case, which
aligns naturally with the Tensor Core tile sizes on the RTX~A6000.

\subsection{Arithmetic intensity in $d$ dimensions}

To understand when the high-dimensional SD-KDE becomes compute-bound, we
estimate FLOPs, bytes moved, and arithmetic intensity for the $d$-dimensional
case.  Let $n_{\text{train}} = k$ and set $n_{\text{test}} = k/8$ as in the
experiments.

\paragraph{Total FLOPs.}
The 16-D implementation consists of three main matrix-multiply stages:
\begin{enumerate}
  \item Score Gram matrix $G = X X^\top$:
        $2 d k^2$ FLOPs.
  \item Score numerator $T = \Phi X$:
        $2 d k^2$ FLOPs, plus $4 k^2$ scalar FLOPs for norms and distance
        terms and $8 k^2$ FLOPs for exponentials (counting each $\exp$ as
        $8$ FLOPs due to the SFU/FP32 ratio).
  \item Final KDE Gram matrix on debiased data:
        $2 d k (k/8)$ FLOPs, plus $4 k (k/8)$ scalar FLOPs and
        $8 k (k/8)$ FLOPs for exponentials.
\end{enumerate}
Aggregating these terms yields
\[
  \begin{aligned}
    \text{FLOPs}_d(k)
    &\approx
    4 d k^2 + 12 k^2 + 2 d \frac{k^2}{8} + 12 \frac{k^2}{8} \\
    &=
    \Bigl(4 d + 12 + \tfrac{d}{4} + \tfrac{3}{2}\Bigr) k^2.
  \end{aligned}
\]
Substituting $d=16$ gives
\[
  \text{FLOPs}_{16}(k)
  \approx
  81.5\,k^2,
\]
which is on the order of $10^{11}$ FLOPs for $k=32\text{k}$.

\paragraph{Bytes moved.}
To better capture implementation details, we count bytes per tile using the
launch parameters that delivered the best runtime ($BLOCK\_M = 64$,
$BLOCK\_N = 1024$).  Each tile loads $64\times d$ query values once
($4\,BLOCK\_M d$ bytes), streams $1024\times d$ training values
($4\,BLOCK\_N d$ bytes), and writes the partial PDF and weighted sums
($4\,BLOCK\_M$ and $4\,BLOCK\_M d$ bytes).  All of this traffic goes to and from
GDDR6, so for $d=16$ a tile moves roughly
\[
  \begin{aligned}
    \text{Bytes}_{\text{tile}}
    &\approx
    4\bigl(BLOCK\_M d + BLOCK\_N d \\
    &\qquad + BLOCK\_M + BLOCK\_M d\bigr) \\
    &= 4\bigl(2\,BLOCK\_M d + BLOCK\_N d \\
    &\qquad + BLOCK\_M\bigr) \\
    &\approx 7.4\times 10^4\ \text{bytes}.
  \end{aligned}
\]
The full problem processes $(k / BLOCK\_M)\times(k / BLOCK\_N)$ such tiles,
so the total GDDR6 traffic is
\[
  \begin{aligned}
    \text{Bytes}_{16}(k)
    &\approx
    \text{Bytes}_{\text{tile}}
    \times
    \frac{k}{BLOCK\_M}
    \times
    \frac{k}{BLOCK\_N} \\
    &\approx
    1.13\,k^2\ \text{bytes}.
  \end{aligned}
\]

\paragraph{Arithmetic intensity.}
Dividing FLOPs by bytes gives
\[
  \begin{aligned}
    I_d(k)
    &=
    \frac{\text{FLOPs}_d(k)}{\text{Bytes}_d(k)} \\
    &\approx
    \frac{\bigl(4 d + 12 + \tfrac{d}{4} + \tfrac{3}{2}\bigr) k^2}
         {4\bigl(\tfrac{9}{8} d k + \tfrac{k}{8}\bigr)} \\
    &\sim
    C(d)\,k
    \quad\text{for large }k,
  \end{aligned}
\]
with
\[
  C(d)
  \approx
  \frac{4 d + 12 + d/4 + 3/2}{4\cdot (9 d/8)}
  =
  \frac{(17/4)\,d + 27/2}{9 d/2}.
\]
For $d=16$ this simplifies to
\[
  I_{16}(k)
  \approx
  \frac{\text{FLOPs}_{16}(k)}{\text{Bytes}_{16}(k)}
  \approx
  \frac{81.5\,k^2}{1.13\,k^2}
  \approx 72\ \text{flops/byte}.
\]
This tile-aware estimate more closely reflects the measured arithmetic
intensity.  Comparing against the A6000 specs (Tensor Core peak of
$155\,\text{TFLOP/s}$ versus $\approx 770\,\text{GB/s}$ of memory bandwidth),
the machine balance is roughly $200$ flops/byte.  Since our kernel sustains
over $70$ flops/byte, it lies well into the compute-bound regime on the
RTX~A6000.  Nsight Compute reports an empirical arithmetic intensity of roughly
$95\,\text{FLOPs/byte}$ for the score kernel on the A6000, which is broadly in
line with the $72\,\text{FLOPs/byte}$ predicted by the tile-level model above.
Minor differences stem from additional bookkeeping work (atomics, reductions)
and from Nsight's finer-grained accounting of texture/cache traffic, but both
figures agree that the kernel operates far above the bandwidth roofline.
Because the kernel mixes Tensor-Core GEMMs with FP32 scalar work (norms,
exponentials, atomics), the effective machine balance sits between the
tensor-core roof ($\approx 200$ flops/byte) and the FP32 roof ($\approx 50$
flops/byte).  The observed $70$--$95$ flops/byte therefore straddle these two
limits: the GEMM portion is partially bandwidth-limited relative to Tensor
Cores, while the scalar portion is compute-limited relative to FP32 ALUs.
